\subsection{Interaction Objective}
\label{text:approach/objective/interactive}
The interaction between robot and pedestrian itself is an abstract quantity and can thus hardly be used for optimization itself. In the following, we consider a scene with the robot facing only one pedestrian to approach this problem, then to generalize the developed concept to interactions with multiple pedestrians.  
\newline
While it is hard to find a measure for the interaction between the robot and the pedestrian in the scene when regarding the problem in general, it can be simplified when a) focussing on one of the interacting agents and b) define the measure to vanish its value. Under these conditions, the problem can be re-formulated as decreasing the disturbance the robot induces on the pedestrian, which is a lot simpler to quantify than the abstract concept of interaction. To do so the un-conditioned\footnote{For brevity the term "un-conditioned" means, not depending on the robot's state. However, the pedestrians trajectories predictions are of course still conditioned on its state history as well as the states and state history of the other pedestrians in the scene, as described in Section \ref{text:approach/formulation}.} pedestrian's trajectory distribution $\distwo[]$ is computed, i.e., the distribution that would occur if no robot would be in the scene, and then compared to its actual conditioned trajectory distribution $\dist[]$.\footnote{When the prediction model requires the input of a robot trajectory, e.g. as an "input format" requirement of a neural network-based model, predicting the un-conditioned trajectory distribution might not be straight-forward. Using a re-trained un-conditioned model is not an option since factors causing a different distribution independent from the robot's trajectory might come into play. Thus, within the project a "pseudo"-robot is used which is located very far away from any pedestrian, hence minimizing its effect on the pedestrian's behavior.} With some general distance measure $\Delta(\cdot, \cdot)$ the interactive objective function given the robot's trajectory $\x_{0:T}$ is defined as:

\begin{equation}
J_{int}^k(\x_{0:T}) = \sum_{t=0}^T \Delta(\dist[k]_t, \distwo[k]_t)
\label{eq:objective_interaction}
\end{equation}

\begin{figure}[!ht]
\begin{center}
\begin{tikzpicture}

    \node[inner sep=0pt] (pedwo) at (-3,1)
    {\includegraphics[width=.05\textwidth]{images/walking.png}};    
    \draw [dotted, ultra thick, name path=A] (pedwo) to[out=180, in=0] node[above] {$\xpedwo[k]$} (-8, 1);
    
    \node[inner sep=0pt] (pedw) at (-3,-1)
    {\includegraphics[width=.05\textwidth]{images/walking.png}};
    \node[inner sep=0pt] (robot) at (-7,-1)
    {\includegraphics[width=.05\textwidth]{images/robot.png}};
    \draw [ultra thick, name path=B] (pedw) to[out=180, in=10] node[below, sloped] {$\xped[k]$}(-8,-3);
    
    \draw[thick, decorate, decoration={brace, amplitude=20pt}] (-1.5,2) -- (-1.5,-2);

    \node[inner sep=0pt] (ped) at (5,0.5)
    {\includegraphics[width=.05\textwidth]{images/walking.png}};
    
    \draw [dotted, ultra thick, name path=A] (ped) to[out=180, in=0] node[above] {$\xpedwo[k]$} (0, 0.5);
    \draw [ultra thick, name path=B] (ped) to[out=180, in=10] node[below, sloped] {$\xped[k]$}(0,-1.5);
        
    \tikzfillbetween[of=A and B]{blue, opacity=0.2};
    \node[] (D) at (1, -0.3){$D_{int}$};
    
\end{tikzpicture}
\end{center}
\caption{Interactive measure in case of deterministic and uni-modal pedestrian trajectory predictions, which trivially is the area enclosed by the un-conditioned (upper) and the conditioned (middle) trajectory.}
\label{img:ado_w_wo_distance_trajectory}
\end{figure}

If the pedestrian's trajectory prediction would be deterministic and unimodal, this measure simply breaks down to the distance between both trajectories, which is the area enclosed by both trajectories in continuous time, as shown in Figure \ref{img:ado_w_wo_distance_trajectory}, or a sum of distance per time-step in discrete time. For probabilistic and multimodal predictions, computing the distance between both distributions is more difficult, especially when demands such as computational cost and differentiability (to be used in an optimization) have to be factored in. 

\begin{figure}[!ht]
\begin{center}
\begin{tikzpicture}

    \node[inner sep=0pt] (ped) at (5,0)
    {\includegraphics[width=.05\textwidth]{images/walking.png}};
    \node[inner sep=0pt] (robot) at (0,0)
    {\includegraphics[width=.05\textwidth]{images/robot.png}};
       
    \draw [dotted, ultra thick, name path=A] (ped) to[out=180, in=0] (0, 2) node[above, sloped] {$\xpedwo[k]$}  to[out=180, in=0] (-4, 0);
    
    \draw [ultra thick, name path=B] (ped) to[out=180, in=10] (0,-2) node[below, sloped] {$\xped[k]$} to[out=180, in=0] (-4, 0);
    
\end{tikzpicture}
\end{center}
\caption{Interactive measure in case of deterministic and uni-modal pedestrian trajectory predictions, which trivially is the area enclosed by the un-conditioned (upper left) and the conditioned (upper right) trajectory.}
\label{img:int_acceleration_reason}
\end{figure}

The correlation between the acceleration carried out on a passenger and its comfort is widely known, e.g., for the driving use case \cite{Hoberock1976}. Therefore, assuming that the same measure applies for correlation between the acceleration the pedestrian itself has to exert (e.g., to evade a dynamic obstacle) and its comfort is likely. Therefore, instead of contrasting the position distributions, it might be valuable to use the velocity or acceleration distributions. Figure \ref{img:int_acceleration_reason} shows a possible scenario in which the robot's presence affects the pedestrian such that the trajectory prediction is mirrored. When the robot is static, and when no other pedestrian is close, both trajectories are equally safe (if the robot is static)  and "comfortable" for the pedestrian and are equal in length, so there is no reason to chose one above the other. However, due to the large enclosed area, a purely position-based distance metric would be non-zero by far, while an acceleration-based distance metric would be zero since the velocities of the pedestrian-only change their sign, not their absolute value.

\subsubsection{Kullback-Leibler Divergence}
A commonly used metric for expressing the distance between two distributions is the Kullback-Leibler Divergence $D_{KL}$, which determines the distance between  some distribution $q$ and another distribution $p$ as:

\begin{equation}
D_{KL} = \int_x q(x) log \frac{q(x)}{p(x)} dx    
\end{equation}

$D_{KL}$ is a well-defined loss term and commonly used in many applications, such as generative deep learning models \cite{Goodfellow2014}\cite{Salzmann2020} (similarly the Jenson-Shannon Divergence). However, it is not analytically defined for some "complex" distributions such as \ac{GMM}, the output distribution of Trajectron \cite{Ivanovic2018}. Methods to approximate the KL-Divergence for \ac{GMM}s have been discussed in \cite{Cui2015} and embrace Monte Carlo sampling, signature quadratic form distance \cite{Beecks2011}, and several more. All of these methods are not computationally feasible for an online application, especially when gradients have to be computed. Other methods simplify the real \ac{GMM} to a single Gaussian, by a weighted average over its parameters, which is not guaranteed to be a meaningful distribution and loses the advantages of predicting multimodal distributions in the first place.
\newline
Regarding the Trajectron \cite{Ivanovic2018} as a prediction model, some intermediate distribution might be used for comparison instead of the output distribution, such as the categorical distribution in its latent space\footnote{In fact, the Trajectron's latent space could not be used as a basis for the interactive objective function anyway, since it does not depend on the robot's trajectory. However, it might be used to assess the similarity between scenarios, which will be discussed in Section \ref{text:approach/runtime/warm_starting}.}. Although this approach might give rise to using standard distance measures such as $D_{KL}$, it would impede tractability and generality of the overall framework, as it would have to be redefined for every other prediction model.

\subsubsection{Trajectory Projection}
Combining computational efficiency and the capability of representing a measure based on the full distribution is hard, as demonstrated in the examples above. However, by exploiting that the predicted distribution $\xped[]_t \sim \dist[]_t$ is a continuous, well-defined distribution with infinite support \footnote{Although not all distributions have infinite support, these properties surely hold for the most commonly used ones in the area of pedestrian prediction such as Gaussians, \ac{GMM}s \cite{Salzmann2020} or non-closed form distribution such as SGAN \cite{Gupta2018}.}, \ref{eq:objective_interaction} can be redefined as the probability of the unconditioned distribution $\distwo[]_t$ with respect to the conditioned distribution $\dist[]_t$ which is equivalent to the integral over the product distribution $\int \int \distwo[]_t \cdot \dist[]_t \, dxdy$.
\newline
For two \ac{GMM}s, the distribution product is not analytically defined though and would involve numerically solving an \ac{ODE} \cite{Schrempf2005} or multi-scale sampling \cite{Ihler2003}. Therefore as simplification, not the full un-conditioned distribution is taken into account but only its mean value $\mathbb{E}[\distwo[]_t]$ and weighted by the conditioned mode importance vector:\footnote{A derivation of the exact objective formulation for a \ac{GMM} as underlying distribution can be found in the appendix.}

\begin{equation}
J_{int}^k = - \sum_{t=0}^T \mathbb{E}_{\xped[k] \sim \dist[k]} \log p(\, \mathbb{E}[\distwo[k]_t] \, | \, \xped[k]_t, \x_t)
\label{eq:objective_interact_prob}
\end{equation}

To deal with reasonable large values, compared to the other objective functions, and for the independence of gradients (sum not product), instead of the \ac{PDF} $p = pdf(x, y)$, its logarithmic value $\log p$ probability is used. Since the product probability should be maximized while the optimization stated in Problem \ref{problem:general} seeks the minimum, the expectation value's negative value is used.
\newline
Equation \ref{eq:objective_interact_prob} is efficient to compute as it can be batched over the full length of the trajectory and all modes. Also, it uses the full distribution and has a unique global minimum when the means of both distributions are identical (at least for Gaussian-like distributions such as \ac{GMM}s). In fact, \ref{eq:objective_interact_prob} is similar to the \ac{ELBO} loss of the Trajectron loss function (compare Equation \ref{eq:trajectron_loss}), which shows that the term is suitable for general optimization, especially concerning the Trajectron model itself. However, $J_{int}$ is generally applicable to all prediction models, that output a probabilistic distribution, independent from multi-modality.
